\section{Lecture 1, Part 0: Data \& Descriptive Statistics}

\begin{frame}
\vfill
\centering
{\color{RLBlue}\large Lecture 1}

\vspace{0.3em}
{\color{RLBlue}\Large\bfseries Part 0}

\vspace{0.6em}
{\Large Data \& Descriptive Statistics}
\vfill
\end{frame}

\begin{frame}{Learning Goals}
\begin{itemize}
  \item Distinguish data types (categorical vs. numerical, discrete vs. continuous) and read frequency distributions, histograms, and box plots.
  \pause
  \item Compute and interpret measures of center (mean, median, mode) and spread (variance, standard deviation, IQR) -- and explain \emph{why} each one is needed, not just how to compute it.
  \pause
  \item Recognize when a single summary statistic hides the real structure of the data (Anscombe's quartet).
\end{itemize}
\end{frame}

\begin{frame}{Data Types}
\begin{columns}[T]
\column{0.48\textwidth}
\textbf{Categorical (qualitative)}
\begin{itemize}
  \item nominal: no natural order (e.g. eye colour, city of birth)
  \pause
  \item ordinal: ordered, unequal spacing (e.g. rating: low/medium/high)
\end{itemize}

\pause
\column{0.48\textwidth}
\textbf{Numerical (quantitative)}
\begin{itemize}
  \item discrete: countable (e.g. number of children in a family)
  \pause
  \item continuous: measured on a scale (e.g. salary, height)
\end{itemize}
\end{columns}

\pause
\begin{rlintuition}
The data type determines which summary statistics and plots are valid: means and standard deviations only make sense for numerical data; categorical data is summarized with counts and proportions.
\end{rlintuition}
\end{frame}

\begin{frame}{Population, Sample, and Notation}
\begin{rldefinition}
A \textbf{population} is every unit we care about; a \textbf{sample} is the subset we actually measure. A sample of size $n$ is written
\[
x_1, x_2, \dots, x_n,
\qquad
x_i = \text{the } i\text{-th observation}.
\]
\end{rldefinition}
\pause
\begin{itemize}
  \item Our 600 salaries are a \textbf{sample}; the population is every employee in the industry -- too large to measure.
  \pause
  \item \textbf{Example:} five salaries (\$1000s) $\{31, 45, 47, 52, 88\}$ give $n = 5$, $x_1 = 31$, $x_3 = 47$, $x_5 = 88$.
  \pause
  \item \textbf{Summation:} $\sum_{i=1}^{n} x_i = x_1 + x_2 + \dots + x_n$ -- ``add up all $n$ observations''.
  \pause
  \item Latin letters ($\bar{x}$, $s$) are computed \emph{from the sample}; Greek letters ($\mu$, $\sigma$) describe the \emph{population} we cannot see.
\end{itemize}
\end{frame}

% ─── Frequency distributions & histograms ───────────────────────────────────
\begin{frame}{Frequency Distributions \& Histograms}
\begin{columns}[T]
\column{0.50\textwidth}
\small
\begin{itemize}
  \item A \textbf{frequency distribution} counts observations per value or class.
  \pause
  \item A \textbf{histogram} plots those counts as bars -- it shows the \emph{shape} of the data: symmetric, skewed, multi-modal.
  \pause
  \item \textbf{Real event:} 600 employee salaries are \emph{right-skewed} -- most sit between \$40k and \$60k, and a thin tail of high earners reaches \$151k.
  \pause
  \item \textbf{Why:} ``the average salary is \$52k'' hides that half the staff earn below \$48k. The shape shows it; the number alone does not.
\end{itemize}
\column{0.48\textwidth}
\centering
\includegraphics[width=\linewidth]{figures/data_probability/data_probability_salary_histogram.pdf}
\end{columns}
\end{frame}

\begin{frame}{Building a Histogram for a Continuous Variable}
\small
A continuous variable almost never repeats a value, so counting each value separately is useless -- group the values into \textbf{classes} first.
\begin{columns}[T]
\column{0.54\textwidth}
\begin{enumerate}
  \item Range $=\max-\min=\$151\text{k}-\$27\text{k}=\$124\text{k}$.
  \pause
  \item Choose the number of classes $k$ (usually 5--10); here $k=7$.
  \pause
  \item Class width $=\text{range}/k=124/7\approx17.7$, rounded up to a convenient \$20k.
  \pause
  \item Count the observations falling in each class.
  \pause
  \item Draw one bar per class, \textbf{bars touching}: the scale is continuous, so gaps would suggest impossible salaries.
\end{enumerate}
\column{0.44\textwidth}
\centering
\scriptsize
\begin{tabular}{lrr}
\hline
Salary class & Freq. & \% \\
\hline
\$20--40k   & 156 & 26.0 \\
\$40--60k   & 311 & 51.8 \\
\$60--80k   &  89 & 14.8 \\
\$80--100k  &  32 &  5.3 \\
\$100--120k &   7 &  1.2 \\
\$120--140k &   3 &  0.5 \\
\$140--160k &   2 &  0.3 \\
\hline
Total       & 600 &      \\
\hline
\end{tabular}
\end{columns}
\end{frame}

% ─── Measures of center: the mean ───────────────────────────────────────────
\begin{frame}{Why Do We Need the Mean?}
\begin{rldefinition}
For a sample $x_1, \dots, x_n$:
\[
\bar{x} = \frac{1}{n}\sum_{i=1}^{n} x_i
\]
\end{rldefinition}
\pause
\begin{itemize}
  \item \textbf{Why:} decisions need one number, not a spreadsheet of 600 salaries. The mean compresses a whole dataset into one value that can be tracked over time or plugged into a formula (expected cost = mean cost $\times$ volume).
  \pause
  \item \textbf{Numerical example:} the five salaries (\$1000s) $\{31, 45, 47, 52, 88\}$ give $\bar{x} = \frac{31+45+47+52+88}{5} = \frac{263}{5} = 52.6$.
  \pause
  \item \textbf{Everyday means:} a student's grade average, the average monthly temperature in a weather report, a car's average fuel consumption, a player's average points per game -- all the same formula.
\end{itemize}
\end{frame}

\begin{frame}{Why Do We Need the Median?}
\begin{rldefinition}
Sort the sample. The \textbf{median} is the middle value: the $\frac{n+1}{2}$-th observation when $n$ is odd, and the average of the two middle observations when $n$ is even.
\end{rldefinition}
\pause
\begin{itemize}
  \item \textbf{Odd $n$:} $\{31, 45, 47, 52, 88\}$ is sorted and $n = 5$, so the median is the 3rd value, $47$.
  \pause
  \item \textbf{Even $n$:} drop the 88 to get $\{31, 45, 47, 52\}$ with $n = 4$; the median is $\frac{45 + 47}{2} = 46$.
  \pause
  \item For the first set $\bar{x} = 52.6$, which is larger than \emph{four of the five} salaries -- the single \$88k earner pulls it up, while the median stays on the middle person.
  \pause
  \item \textbf{Why:} the median describes a \emph{typical} member of the sample; the mean describes what each would get if the total were shared equally.
\end{itemize}
\end{frame}

\begin{frame}{Real Event: When the Mean Misleads}
\begin{columns}[T]
\column{0.50\textwidth}
\small
\begin{itemize}
  \item Nine people are in a bar, net worth \$20K--\$150K. Mean $\approx$ \$97K, median $\approx$ \$121K -- close together.
  \pause
  \item A billionaire walks in. The \textbf{mean} jumps to \$15{,}000{,}087{,}166 -- nobody in the room is close to that. The \textbf{median} barely moves, to \$122K.
  \pause
  \item \textbf{Why:} the median only cares about the middle-ranked value; the mean sums every value, so one huge outlier dominates the sum.
\end{itemize}
\column{0.48\textwidth}
\centering
\includegraphics[width=\linewidth]{figures/data_probability/data_probability_billionaire_bar.pdf}
\end{columns}
\end{frame}

\begin{frame}{Mean vs. Median vs. Mode: Choosing a Center}
\begin{itemize}
  \item \textbf{Why this matters in practice:} national statistics agencies report \emph{median} household income, not mean -- a handful of billionaires would make ``mean income'' meaningless for a typical household. House prices are reported the same way: a few mansions would drag the mean far above what an ordinary buyer pays.
  \pause
  \item The \textbf{mode} is the most frequent value -- the only center defined for categorical data (the most common blood type, the shoe size a shop sells most).
  \pause
  \item \textbf{Rule of thumb:} mean for roughly symmetric data and for totals (a shop's daily takings); median for skewed data or extreme outliers (income, house prices, waiting times); mode for categorical data.
\end{itemize}
\end{frame}

% ─── Measures of spread: variance & std ─────────────────────────────────────
\begin{frame}{Why Do We Need Variance?}
\begin{rldefinition}
\[
s^2 = \frac{1}{n-1}\sum_{i=1}^{n}(x_i - \bar{x})^2
\qquad
s = \sqrt{s^2}
\]
\end{rldefinition}
\pause
\begin{itemize}
  \item \textbf{Why:} the mean alone cannot distinguish ``reliable'' from ``risky.'' Two datasets can share the exact same mean and still behave completely differently.
  \pause
  \item \textbf{Numerical example:} your commute takes $\{30, 29, 31, 30, 30\}$ minutes by bus and $\{18, 42, 25, 35, 30\}$ by car. Both average $30$ minutes, but $s^2_{\text{bus}} = \frac{2}{4} = 0.5$ and $s^2_{\text{car}} = \frac{338}{4} = 84.5$ -- by car you are sometimes 12 minutes late.
  \pause
  \item $s^2$ divides by $n-1$ (Bessel's correction) to give an unbiased estimate of the \emph{population} variance from a sample.
\end{itemize}
\end{frame}

\begin{frame}{Real Event: Same Mean, Different Variance}
\begin{columns}[T]
\column{0.50\textwidth}
\small
\begin{itemize}
  \item \textbf{Bottling line:} the label says $500$ml; the law allows $495$--$505$ml. Two machines both average exactly $500$ml.
  \pause
  \item \textbf{Machine A}, std $=2$ml: the limits sit $2.5$ standard deviations out, so \textbf{11 of 1000} bottles are rejected.
  \pause
  \item \textbf{Machine B}, std $=4$ml: the same limits are only $1.25$ standard deviations out -- \textbf{210 of 1000} rejected, one bottle in five.
  \pause
  \item Doubling the spread multiplied the waste by $19$; the mean never moved.
\end{itemize}
\column{0.48\textwidth}
\centering
\includegraphics[width=\linewidth]{figures/data_probability/data_probability_variance_bottles.pdf}
\end{columns}
\end{frame}

\begin{frame}{Why This Matters: Variance Determines Risk}
\begin{itemize}
  \item The bottler must control the \emph{spread}, not just the average, because the spread is what sets the share of mislabelled bottles. (This is the idea behind ``Six Sigma'' quality control.)
  \pause
  \item Same story elsewhere: a bus line with a reliable timetable, a clinic with predictable waiting times, an oven that holds its temperature -- the average alone never says how often things go wrong.
  \pause
  \item \textbf{Later in this course:} a model whose accuracy swings from run to run is riskier than a steady one with the same average -- the spread decides again.
\end{itemize}
\end{frame}

\begin{frame}{Quartiles and the Five-Number Summary}
\begin{rldefinition}
Sort the data and split it at the median. $Q_1$ is the median of the lower half and $Q_3$ the median of the upper half, so a quarter of the data lies below $Q_1$ and a quarter above $Q_3$. The \textbf{five-number summary} is $\min$, $Q_1$, median, $Q_3$, $\max$.
\end{rldefinition}
\pause
\begin{itemize}
  \item \textbf{Worked example} -- nine salaries (\$1000s), sorted, median in bold:
  \[
    31,\; 38,\; 42,\; 45,\; \mathbf{47},\; 49,\; 52,\; 61,\; 88
  \]
  \pause
  \item Lower half $\{31,38,42,45\}$ gives $Q_1 = \frac{38+42}{2} = 40$; upper half $\{49,52,61,88\}$ gives $Q_3 = \frac{52+61}{2} = 56.5$.
  \pause
  \item Five-number summary: $31,\; 40,\; 47,\; 56.5,\; 88$ -- the centre \emph{and} the spread, in five numbers.
\end{itemize}
\end{frame}

\begin{frame}{IQR \& Box Plots}
\begin{columns}[T]
\column{0.50\textwidth}
\small
\begin{itemize}
  \item \textbf{IQR} $= Q_3 - Q_1$: spread of the middle 50\% -- like the median, robust to extreme outliers.
  \pause
  \item Outliers: beyond $Q_1 - 1.5\,\text{IQR}$ or $Q_3 + 1.5\,\text{IQR}$. For the nine salaries, $\text{IQR} = 16.5$ and the upper fence is $81.25$, so \$88k is flagged.
  \pause
  \item A \textbf{box plot} draws $Q_1$, median, $Q_3$, whiskers, outliers -- faster to compare groups than several histograms.
  \pause
  \item \textbf{Real event:} waiting times at three branches of a service centre. Branch C is both the slowest \emph{and} the most variable -- a single average would not show \emph{which} branch to fix.
\end{itemize}
\column{0.48\textwidth}
\centering
\includegraphics[width=\linewidth]{figures/data_probability/data_probability_boxplot_branches.pdf}
\end{columns}
\end{frame}

\begin{frame}{Why We Always Plot the Data: Anscombe's Quartet}
\centering
\includegraphics[width=0.96\textwidth]{figures/data_probability/data_probability_anscombe.pdf}

\vspace{0.4em}
\footnotesize All four datasets share mean $x=9.0$, mean $y\approx7.50$, variance $x=11.0$, correlation $\approx 0.816$, and the same regression line $y = 3.00 + 0.5x$.
\end{frame}

\begin{frame}{Anscombe's Quartet: The Lesson}
\begin{itemize}
  \item \textbf{Real event:} statistician Francis Anscombe published these four datasets in 1973 specifically to warn against trusting summary statistics without plotting the data.
  \pause
  \item Set I: a normal linear relationship with scatter. Set II: a perfect \emph{curve} -- a line is the wrong model entirely. Set III: a perfect line broken by one outlier that pulls the fitted slope off. Set IV: no relationship at all except one high-leverage point that fakes a strong correlation.
  \pause
  \item \textbf{Why this matters for ML:} reporting ``$R^2=0.67$, mean=9.0, std=3.3'' for a dataset tells you almost nothing about whether a linear model is appropriate, whether outliers are driving your fit, or whether the relationship is even real -- always visualize residuals and scatter plots, not just summary numbers, before trusting a model.
\end{itemize}
\end{frame}
